\section{Related Works}
\label{related}

\textbf{Unified Multimodal Models}\quad
A classic strategy for integrating visual understanding and generation within a single MLLM is to externally connect an LLM with a Diffusion Model~\citep{sun2024Emu2,dong2024dreamllm,metaquery,blip3o}. However, pure AR architectures offer a more elegant, fully end-to-end solution by unifying both tasks within the same autoregressive framework.
Representative works like Chameleon~\citep{yu2023CM3leon,team2024chameleon} and Emu3~\citep{wang2024emu3}, have demonstrated the feasibility of jointly modeling vision and language through a unified next-token prediction objective. Specifically, visual inputs are first tokenized into visual tokens. These visual tokens are then interleaved with text tokens to construct a multimodal sequence.
However, these pure AR architectures introduce generative capabilities at the cost of considerably weaker visual understanding. An empirical explanation for this~\citep{wu2025vilau,xie2025musevl} is that their vision tokenizers are trained solely for reconstruction and thus primarily captures low-level visual details for generation rather than the high-level semantics required for vision–language understanding.
% Therefore, improving the vision–language understanding performance of AR-based models necessitates a vision tokenizer that can effectively capture both low-level details for generation and high-level semantics for understanding.

A straightforward way to bypass such a conflict is to employ two heterogeneous vision encoders~\citep{wu2024janus,chen2025januspro,deng2025bagel}: a semantic tokenizer (\eg CLIP) for understanding and a reconstruction-based tokenizer (\eg VQ-VAE) for generation. Yet this design inevitably adds extra modules and structural complexity, making understanding and generation two loosely coupled systems with distinct pathways rather than a truly unified model.
In contrast, the text modality relies on a single tokenizer (\eg, BPE)~\citep{sennrich2015bpe} that discretizes text into a unified token space. This ensures a consistent input–output space: the input tokens that provide signals for understanding and the output tokens produced during generation share the same vocabulary.
This unified design allows LLMs to seamlessly integrate text understanding and generation within the next-token prediction paradigm, thereby supporting broad generalization across diverse linguistic tasks.
Therefore, the visual modality urgently requires a tokenizer that, like text tokenizers, can support both understanding and generation within a unified, coherent token space.

\textbf{Unified Visual Tokenizers}\quad
Recent research has actively explored solutions in this direction. VILA-U~\citep{wu2025vilau} and MUSE-VL~\citep{xie2025musevl} strive to build a unified tokenizer by jointly training on both reconstruction and semantic objectives.
However, due to the inherent disparity between semantic and texture features, they struggle to strike an optimal balance between the two objectives, resulting in subpar performance in both tasks.
As discussed in FQGAN~\citep{bai2024fqgan}, decomposing the codebook in a divide-and-conquer manner may offer a more fundamental solution to this conflict.
TokenFlow~\citep{qu2024tokenflow} employs separate codebooks with a shared-mapping mechanism. However, key differences set our approach apart: (i) TokenFlow relies on distinct vision towers to extract semantic and low-level features, rather than leveraging a unified architecture; (ii) the shared IDs obtained through the shared-mapping mechanism may not be the optimal matches for either semantics or texture, potentially introducing additional losses in both domains.
